\chapter{2010年大纲II卷（理）}

\section{单选题}

\begin{problem}
复数 \(\left(\frac{3 - \i}{1 + \i}\right)^2 =\)（\quad）

\choices
    {\(-3 - 4\i\)}
    {\(-3 + 4\i\)}
    {\(3 - 4\i\)}
    {\(3 + 4\i\)}
\end{problem}

\begin{answer}
A
\end{answer}

\begin{solution}
先化简分式，得
\[
\frac{3-\i}{1+\i}=\frac{(3-\i)(1-\i)}{(1+\i)(1-\i)}=\frac{2-4\i}{2}=1-2\i
\]
因此
\[
\left(\frac{3-\i}{1+\i}\right)^2=(1-2\i)^2=1-4\i+4\i^2=-3-4\i
\]
所以选 A.
\end{solution}

\begin{problem}
函数 \( y = \frac{1 + \ln(x - 1)}{2} \)（\(x > 1\)）的反函数是（\quad）

\choices
    {\( y = \e^{2x + 1} - 1 \)（\(x > 0\)）}
    {\( y = \e^{2x - 1} + 1 \)（\(x > 0\)）}
    {\( y = \e^{2x + 1} - 1 \)（\(x \in \R\)）}
    {\( y = \e^{2x - 1} + 1 \)（\(x \in \R\)）}
\end{problem}

\begin{answer}
D
\end{answer}

\begin{solution}
设原函数中的自变量为 \(x\)，函数值为 \(y\)，则
\[
2y=1+\ln(x-1)
\]
所以
\(\ln(x-1)=2y-1\)，\(x-1=\e^{2y-1}\)，\(x=\e^{2y-1}+1\).
原函数的值域为 \(\R\)，交换 \(x\)，\(y\) 后，反函数为
\(y=\e^{2x-1}+1\)，\((x\in\R)\).
所以选 D.
\end{solution}

\begin{problem}
若变量 \(x\)，\(y\) 满足约束条件 \(\begin{cases}
x \geqslant -1,\\
y \geqslant x,\\
3x + 2y \leqslant 5,
\end{cases}\) 则 \(z = 2x + y\) 的最大值为（\quad）

\choices
    {\(1\)}
    {\(2\)}
    {\(3\)}
    {\(4\)}
\end{problem}

\begin{answer}
C
\end{answer}

\begin{solution}
由 \(y\geqslant x\) 和 \(3x+2y\leqslant5\)，可得
\[
x\leqslant y\leqslant\frac{5-3x}{2}
\]
从而 \(x\leqslant1\). 结合 \(x\geqslant-1\)，有 \(-1\leqslant x\leqslant1\). 又因为
\[
z=2x+y\leqslant2x+\frac{5-3x}{2}=\frac{x+5}{2}\leqslant3
\]
当 \(x=1\)，\(y=1\) 时，三个约束条件均成立，且 \(z=3\)，故最大值为 \(3\). 所以选 C.
\end{solution}

\begin{problem}
如果等差数列 \(\{a_{n}\}\) 中，\(a_{3} + a_{4} + a_{5} = 12\)，那么 \(a_{1} + a_{2} + \dots + a_{7} =\)（\quad）

\choices
    {\(14\)}
    {\(21\)}
    {\(28\)}
    {\(35\)}
\end{problem}

\begin{answer}
C
\end{answer}

\begin{solution}
等差数列满足中间项等于相邻三项的平均值，因此
\(a_3+a_4+a_5=3a_4=12\)，\(a_4=4\).
又 \(a_1+a_7=a_2+a_6=a_3+a_5=2a_4\)，所以
\[
a_1+a_2+\cdots+a_7=7a_4=28
\]
所以选 C.
\end{solution}

\begin{problem}
不等式 \(\frac{x^2 - x - 6}{x - 1} > 0\) 的解集为（\quad）

\choices
    {\(\{x\mid x < -2 \text{或} x > 3\}\)}
    {\(\{x\mid x < -2 \text{或} 1 < x < 3\}\)}
    {\(\{x\mid -2 < x < 1 \text{或} x > 3\}\)}
    {\(\{x\mid -2 < x < 1 \text{或} 1 < x < 3\}\)}
\end{problem}

\begin{answer}
C
\end{answer}

\begin{solution}
将分子因式分解，得
\[
\frac{x^2-x-6}{x-1}=\frac{(x+2)(x-3)}{x-1}
\]
临界点为 \(-2\)，\(1\)，\(3\)，其中 \(x=1\) 不在定义域内. 在区间
\((-\infty,-2)\)，\((-2,1)\)，\((1,3)\)，\((3,+\infty)\)
上分别判断三个因式的符号，原式的符号依次为负、正、负、正. 因为要求原式大于 \(0\)，且 \(-2\)，\(3\) 也不能取，所以解集为
\[
\{x\mid -2<x<1\text{或}x>3\}
\]
所以选 C.
\end{solution}

\begin{problem}
将标号为 \(1\)，\(2\)，\(3\)，\(4\)，\(5\)，\(6\) 的 \(6\) 张卡片放入 \(3\) 个不同的信封中，若每个信封放 \(2\) 张，其中标号为 \(1\)，\(2\) 的卡片放入同一信封，则不同的放法共有（\quad）

\choices
    {\(12\text{ 种}\)}
    {\(18\text{ 种}\)}
    {\(36\text{ 种}\)}
    {\(54\text{ 种}\)}
\end{problem}

\begin{answer}
B
\end{answer}

\begin{solution}
先把标号为 \(1\)，\(2\) 的卡片放入同一个信封，有 \(3\) 种选择. 选定该信封后，剩下的 \(4\) 张卡片须分别放入另外两个有区别的信封，指定其中一个信封放入 \(2\) 张卡片，另一个信封自动得到其余 \(2\) 张，放法有 \(\mathrm C_4^2=6\) 种. 因此总放法数为
\[
3\mathrm C_4^2=3\times6=18
\]
所以选 B.
\end{solution}

\begin{problem}
为了得到函数 \( y = \sin \left(2x - \frac{\pi}{3}\right) \) 的图象，只需把函数 \( y = \sin \left(2x + \frac{\pi}{6}\right) \) 的图象（\quad）

\choices
    {向左平移 \(\frac{\pi}{4}\text{ 个长度单位}\)}
    {向右平移 \(\frac{\pi}{4}\text{ 个长度单位}\)}
    {向左平移 \(\frac{\pi}{2}\text{ 个长度单位}\)}
    {向右平移 \(\frac{\pi}{2}\text{ 个长度单位}\)}
\end{problem}

\begin{answer}
B
\end{answer}

\begin{solution}
把原函数图象向右平移 \(h\) 个长度单位后，函数变为
\[
y=\sin\left(2(x-h)+\frac{\pi}{6}\right)=\sin\left(2x-2h+\frac{\pi}{6}\right)
\]
要使它等于 \(y=\sin(2x-\frac{\pi}{3})\)，应满足
\(-2h+\frac{\pi}{6}=-\frac{\pi}{3}\)，\(h=\frac{\pi}{4}\).
所以应向右平移 \(\frac{\pi}{4}\) 个长度单位，选 B.
\end{solution}

\begin{problem}
\(\triangle ABC\) 中，点 \(D\) 在边 \(AB\) 上，\(CD\) 平分 \(\angle ACB\). 若 \(\overrightarrow{CB} = \bs{a}\)，\(\overrightarrow{CA} = \bs{b}\)，\(|\bs{a}| = 1\)，\(|\bs{b}| = 2\)，则 \(\overrightarrow{CD} =\)（\quad）

\choices
    {\(\frac{1}{3} \bs{a} + \frac{2}{3} \bs{b}\)}
    {\(\frac{2}{3} \bs{a} + \frac{1}{3} \bs{b}\)}
    {\(\frac{3}{5} \bs{a} + \frac{4}{5} \bs{b}\)}
    {\(\frac{4}{5} \bs{a} + \frac{3}{5} \bs{b}\)}
\end{problem}

\begin{answer}
B
\end{answer}

\begin{solution}
由角平分线定理，
\[
\frac{AD}{DB}=\frac{AC}{CB}=\frac{|\bs b|}{|\bs a|}=2
\]
因此 \(AD:DB=2:1\)，点 \(D\) 将从 \(A\) 到 \(B\) 的向量按比例分点，故
\[
\overrightarrow{CD}=\overrightarrow{CA}+\frac{2}{3}\overrightarrow{AB}
=\bs b+\frac{2}{3}(\bs a-\bs b)
=\frac{2}{3}\bs a+\frac{1}{3}\bs b
\]
所以选 B.
\end{solution}

\begin{problem}
已知正四棱锥 \(S - ABCD\) 中，\(SA = 2\sqrt{3}\)，那么当该棱锥的体积最大时，它的高为（\quad）

\choices
    {\(1\)}
    {\(\sqrt{3}\)}
    {\(2\)}
    {\(3\)}
\end{problem}

\begin{answer}
C
\end{answer}

\begin{solution}
设正方形底面的边长为 \(s\)，棱锥的高为 \(h\)，底面中心为 \(O\). 因为
\(OA=\frac{s}{\sqrt{2}}\)，\(SA^2=h^2+OA^2=12\)
所以
\[
s^2=2(12-h^2)
\]
棱锥体积为
\[
V=\frac{1}{3}s^2h=\frac{2}{3}(12-h^2)h
\]
令 \(u=h^2\)，则 \(0<u<12\)，且
\(V=\frac{2}{3}(12-u)\sqrt{u}\)，\(V'(u)=\frac{4-u}{\sqrt{u}}\).
因此 \(V\) 在 \(0<u<4\) 时递增，在 \(4<u<12\) 时递减，体积最大时 \(u=4\)，从而 \(h=2\). 所以选 C.
\end{solution}

\begin{problem}
若曲线 \( y = x^{-\frac{1}{2}} \) 在点 \( \left(a, a^{-\frac{1}{2}}\right) \) 处的切线与两个坐标轴围成的三角形的面积为 \(18\)，则 \(a =\)（\quad）

\choices
    {\(64\)}
    {\(32\)}
    {\(16\)}
    {\(8\)}
\end{problem}

\begin{answer}
A
\end{answer}

\begin{solution}
函数 \(y=x^{-1/2}\) 的导数为
\[
y'=-\frac{1}{2}x^{-3/2}
\]
因此曲线在 \(\left(a,a^{-1/2}\right)\) 处的切线斜率为 \(-1/(2a^{3/2})\)，切线方程为
\[
y-a^{-1/2}=-\frac{1}{2a^{3/2}}(x-a)
\]
令 \(y=0\)，得切线在 \(x\) 轴上的截距为 \(3a\)；令 \(x=0\)，得切线在 \(y\) 轴上的截距为 \(3/(2\sqrt a)\). 因而三角形面积为
\[
\frac{1}{2}\cdot3a\cdot\frac{3}{2\sqrt a}=\frac{9\sqrt a}{4}=18
\]
所以 \(\sqrt a=8\)，\(a=64\)，选 A.
\end{solution}

\begin{problem}
与正方体 \(ABCD - A_{1}B_{1}C_{1}D_{1}\) 的三条棱 \(AB\)，\(CC_{1}\)，\(A_{1}D_{1}\) 所在直线的距离相等的点（\quad）

\choices
    {有且只有 \(1\text{ 个}\)}
    {有且只有 \(2\text{ 个}\)}
    {有且只有 \(3\text{ 个}\)}
    {有无数个}
\end{problem}

\begin{answer}
D
\end{answer}

\begin{solution}
设正方体边长为 \(1\)，建立直角坐标系：
\(A(0,0,0)\)，\(B(1,0,0)\)，\(C(1,1,0)\)，\(A_1(0,0,1)\)，\(C_1(1,1,1)\)，\(D_1(0,1,1)\).
设空间中任一点为 \(P(x,y,z)\). 则
\(d^2(P,AB)=y^2+z^2\)，\(d^2(P,CC_1)=(x-1)^2+(y-1)^2\)
\[
d^2(P,A_1D_1)=x^2+(z-1)^2
\]
注意直线 \(A_1D_1\) 平行于 \(y\) 轴，因此到它的距离由 \(x\)，\(z\) 两个坐标决定.

对任意 \(t\in\R\)，取
\[
P_t=(t,1-t,t)
\]
则
\(d^2(P_t,AB)=(1-t)^2+t^2\)，\(d^2(P_t,CC_1)=(t-1)^2+t^2\)，\(d^2(P_t,A_1D_1)=t^2+(t-1)^2\).
三者相等，故每个 \(P_t\) 都满足条件. 由于 \(t\) 可以取任意实数，满足条件的点有无数个，所以选 D.
\end{solution}

\begin{problem}
已知椭圆 \(C: \frac{x^2}{a^2} + \frac{y^2}{b^2} = 1\)（\(a > b > 0\)）的离心率为 \(\frac{\sqrt{3}}{2}\)，过右焦点 \(F\) 且斜率为 \(k\)（\(k > 0\)）的直线与 \(C\) 相交于 \(A\)，\(B\) 两点. 若 \(\overrightarrow{AF} = 3\overrightarrow{FB}\)，则 \(k =\)（\quad）

\choices
    {\(1\)}
    {\(\sqrt{2}\)}
    {\(\sqrt{3}\)}
    {\(2\)}
\end{problem}

\begin{answer}
B
\end{answer}

\begin{solution}
设椭圆焦距的一半为 \(c\). 由离心率得
\(\frac{c}{a}=\frac{\sqrt3}{2}\)，\(c^2=\frac{3a^2}{4}\)，\(b^2=a^2-c^2=\frac{a^2}{4}\).
取右焦点 \(F=(c,0)\)，并设
\[
\overrightarrow{FB}=(t,kt)
\]
由 \(\overrightarrow{AF}=3\overrightarrow{FB}\)，有
\(B=F+(t,kt)\)，\(A=F-3(t,kt)\).
直线 \(AB\) 上的点可写成 \(F+s(1,k)\). 将其代入椭圆方程并乘以 \(a^2\)，得
\[
(c+s)^2+4k^2s^2=a^2
\]
即
\[
(1+4k^2)s^2+2cs-\frac{a^2}{4}=0
\]
该方程的两个根为 \(-3t\)，\(t\)，由韦达定理，
\(-2t=-\frac{2c}{1+4k^2}\)，\(-3t^2=-\frac{a^2}{4(1+4k^2)}\).
由第一式 \(t=c/(1+4k^2)\)，代入第二式并利用 \(c^2=3a^2/4\)，得
\[
\frac{9}{(1+4k^2)^2}=\frac{1}{1+4k^2}
\]
因为 \(1+4k^2>0\)，所以 \(1+4k^2=9\)，从而 \(k^2=2\). 又 \(k>0\)，故 \(k=\sqrt2\). 所以选 B.
\end{solution}

\section{填空题}

\begin{problem}
已知 \(\alpha\) 是第二象限的角，\(\tan\left(\pi + 2\alpha\right) = -\frac{4}{3}\)，则 \(\tan\alpha =\)\fillinblank{}.
\end{problem}

\begin{answer}
\(-\frac{1}{2}\)
\end{answer}

\begin{solution}
由正切函数的周期性，
\[
\tan(\pi+2\alpha)=\tan2\alpha=-\frac43
\]
令 \(t=\tan\alpha\)，则
\[
\frac{2t}{1-t^2}=-\frac43
\]
因为题目中的正切有定义，所以 \(t^2\ne1\)，交叉相乘得
\(6t=-4(1-t^2)\)，\(2t^2-3t-2=0\)
即
\[
(2t+1)(t-2)=0
\]
所以 \(t=-\frac12\) 或 \(t=2\). 由于 \(\alpha\) 是第二象限的角，\(\tan\alpha<0\)，故 \(\tan\alpha=-\frac12\).
\end{solution}

\begin{problem}
若 \(\left(x - \frac{a}{x}\right)^9\) 的展开式中 \(x^3\) 的系数是\(-84\)，则 \(a =\)\fillinblank{}.
\end{problem}

\begin{answer}
\(1\)
\end{answer}

\begin{solution}
展开式通项中取 \(-a/x\) 的 \(k\) 次幂时，项为
\[
\mathrm C_9^k x^{9-k}\left(-\frac{a}{x}\right)^k
=\mathrm C_9^k(-a)^k x^{9-2k}
\]
要得到 \(x^3\) 项，必须满足
\(9-2k=3\)，\(k=3\).
因此 \(x^3\) 的系数为
\[
\mathrm C_9^3(-a)^3=-84a^3
\]
由题意 \(-84a^3=-84\)，得 \(a^3=1\)，所以 \(a=1\).
\end{solution}

\begin{problem}
已知抛物线 \(C: y^2 = 2px\)（\(p > 0\)）的准线为 \(l\)，过 \(M(1,0)\) 且斜率为 \(\sqrt{3}\) 的直线与 \(l\) 相交于点 \(A\)，与 \(C\) 的一个交点为 \(B\). 若 \(\overrightarrow{AM} = \overrightarrow{MB}\)，则 \(p =\fillinblank{}\).
\end{problem}

\begin{answer}
\(2\)
\end{answer}

\begin{solution}
抛物线 \(y^2=2px\) 的准线为 \(x=-p/2\). 过 \(M(1,0)\) 且斜率为 \(\sqrt3\) 的直线为
\[
y=\sqrt3(x-1)
\]
因此
\[
A=\left(-\frac p2,-\sqrt3\left(\frac p2+1\right)\right)
\]
由 \(\overrightarrow{AM}=\overrightarrow{MB}\)，有 \(B=2M-A\)，所以
\[
B=\left(2+\frac p2,\sqrt3\left(\frac p2+1\right)\right)
\]
将 \(B\) 代入抛物线方程，得
\[
3\left(\frac p2+1\right)^2=2p\left(2+\frac p2\right)=4p+p^2
\]
整理得
\(\frac{p^2}{4}+p-3=0\)，\((p+6)(p-2)=0\).
因为 \(p>0\)，所以 \(p=2\).
\end{solution}

\begin{problem}
已知球 \(O\) 的半径为 4，圆 \(M\) 与圆 \(N\) 为该球的两个小圆，\(AB\) 为圆 \(M\) 与圆 \(N\) 的公共弦，\(AB = 4\). 若 \(OM = ON = 3\)，则两圆圆心的距离 \(MN = \)\fillinblank{}.
\end{problem}

\begin{answer}
\(3\)
\end{answer}

\begin{solution}
设 \(H\) 为公共弦 \(AB\) 的中点. 由于 \(OH\perp AB\)，球半径为 \(4\)，且 \(AB=4\)，有
\[
OH=\sqrt{4^2-\left(\frac{AB}{2}\right)^2}=\sqrt{16-4}=2\sqrt3
\]
对圆 \(M\)，其圆心 \(M\) 到弦 \(AB\) 垂直，且小圆半径满足
\[
MA=\sqrt{4^2-OM^2}=\sqrt7
\]
所以
\[
MH=\sqrt{MA^2-\left(\frac{AB}{2}\right)^2}=\sqrt{7-4}=\sqrt3
\]
同理 \(NH=\sqrt3\). 又 \(OM\perp MH\)，\(ON\perp NH\)，故在直角三角形 \(OMH\) 中
\(\cos\angle MHO=\frac{MH}{OH}=\frac12\)，\(\angle MHO=60^\circ\).
由于 \(MH\perp AB\)，\(NH\perp AB\)，点 \(M\)，\(N\) 都位于过 \(H\) 且垂直于 \(AB\) 的平面内. 在这个平面内，满足与 \(HO\) 成 \(60^\circ\) 角且到 \(H\) 的距离为 \(\sqrt3\) 的位置只有 \(HO\) 两侧的两个点；圆 \(M\) 与圆 \(N\) 不同，所以 \(M\)，\(N\) 位于 \(HO\) 的两侧，从而
\[
\angle MHN=120^\circ
\]
在三角形 \(MHN\) 中，应用余弦定理得
\[
MN^2=MH^2+NH^2-2MH\cdot NH\cos120^\circ=3+3+3=9
\]
所以 \(MN=3\).
\end{solution}

\section{解答题}

\begin{problem}
\(\triangle ABC\) 中，\(D\) 为边 \(BC\) 上的一点，\(BD = 33\)，\(\sin B = \frac{5}{13}\)，\(\cos \angle ADC = \frac{3}{5}\)，求 \(AD\).
\end{problem}

\begin{solution}
设 \(\angle ADC=\theta\)，则 \(\cos\theta=\frac{3}{5}\)，且 \(0<\theta<\frac{\pi}{2}\)，所以 \(\sin\theta=\frac{4}{5}\). 由于 \(B\)，\(D\)，\(C\) 三点共线，\(\angle ADB=\pi-\theta\)，从而 \(\cos\angle ADB=-\frac{3}{5}\)，\(\sin\angle ADB=\frac{4}{5}\). 又因为 \(\angle ADB>\frac{\pi}{2}\)，在三角形 \(ABD\) 中有 \(B<\theta<\frac{\pi}{2}\)，故 \(\cos B=\frac{12}{13}\).

由三角形内角和及正弦的差角公式，得
\[
\sin\angle BAD=\sin(\theta-B)=\sin\theta\cos B-\cos\theta\sin B=\frac{4}{5}\cdot\frac{12}{13}-\frac{3}{5}\cdot\frac{5}{13}=\frac{33}{65}
\]

在三角形 \(ABD\) 中，由正弦定理得
\[
\frac{AD}{\sin B}=\frac{BD}{\sin\angle BAD}
\]
所以 \(AD=33\cdot\frac{5}{13}\cdot\frac{65}{33}=25\)，即 \(AD=25\).
\end{solution}

\begin{problem}
已知数列 \(\{a_{n}\}\) 的前 \(n\) 项和 \(S_{n} = (n^{2} + n)\cdot 3^{n}\).
\begin{enumerate}
\item 求 \(\displaystyle\lim_{n\to\infty}\frac{a_{n}}{S_{n}}\)；
\item 证明：\(\frac{a_{1}}{1^{2}} + \frac{a_{2}}{2^{2}} + \cdots + \frac{a_{n}}{n^{2}} > 3^{n}\).
\end{enumerate}
\end{problem}

\begin{solution}
由 \(a_n=S_n-S_{n-1}\) 及 \(S_0=0\)，得
\[
a_n=(n^2+n)3^n-\bigl((n-1)^2+(n-1)\bigr)3^{n-1}=2n(n+2)3^{n-1}
\]

\begin{enumerate}
\item 因此
\[
\frac{a_n}{S_n}=\frac{2n(n+2)3^{n-1}}{n(n+1)3^n}=\frac{2(n+2)}{3(n+1)}
\]
从而 \(\displaystyle\lim_{n\to\infty}\frac{a_n}{S_n}=\frac{2}{3}\).

\item 对任意正整数 \(k\)，有
\[
\frac{a_k}{k^2}=\frac{2(k+2)3^{k-1}}{k}=2\cdot 3^{k-1}+\frac{4\cdot 3^{k-1}}{k}>2\cdot 3^{k-1}
\]
由等比数列求和公式可得
\[
\sum_{k=1}^{n}\frac{a_k}{k^2}
=2\sum_{k=1}^{n}3^{k-1}+4\sum_{k=1}^{n}\frac{3^{k-1}}{k}
=3^n-1+4\sum_{k=1}^{n}\frac{3^{k-1}}{k}>3^n
\]
其中最后一个不等式由 \(\sum_{k=1}^{n}\frac{3^{k-1}}{k}\geqslant1\) 得出. 因此
\[
\frac{a_1}{1^2}+\frac{a_2}{2^2}+\cdots+\frac{a_n}{n^2}>3^n
\]
\end{enumerate}
\end{solution}

\begin{problem}
如图，直三棱柱 \(ABC - A_{1}B_{1}C_{1}\) 中，\(AC = BC\)，\(AA_{1} = AB\)，\(D\) 为 \(BB_{1}\) 的中点，\(E\) 为 \(AB_{1}\) 上的一点，\(AE = 3EB_{1}\).
\begin{enumerate}
\item 证明：\(DE\) 为异面直线 \(AB_{1}\) 与 \(CD\) 的公垂线.
\item 设异面直线 \(AB_{1}\) 与 \(CD\) 的夹角为 \(45^{\circ}\)，求二面角 \(A_{1} - AC_{1} - B_{1}\) 的大小.
\end{enumerate}
\bitmapfigure[width=0.44\linewidth]{img/2010/syllabus_paper_2_science/q19_fig1.png}
\end{problem}

\begin{solution}
设 \(AB=h>0\)，取直角坐标系，使 \(AB\) 的中点为原点，\(AB\) 为 \(x\) 轴，棱柱的侧棱方向为 \(z\) 轴. 由 \(AC=BC\)，可设 \(C=(0,q,0)\)，其中 \(q>0\). 于是 \(A=\left(-\frac{h}{2},0,0\right)\)，\(B=\left(\frac{h}{2},0,0\right)\)，\(C_1=(0,q,h)\)，\(A_1=\left(-\frac{h}{2},0,h\right)\)，\(B_1=\left(\frac{h}{2},0,h\right)\). 又因 \(D\) 为 \(BB_1\) 的中点，且 \(AE=3EB_1\)，得 \(D=\left(\frac{h}{2},0,\frac{h}{2}\right)\)，\(E=\left(\frac{h}{4},0,\frac{3h}{4}\right)\).

\begin{enumerate}
\item 于是 \(\overrightarrow{AB_1}=(h,0,h)\)，\(\overrightarrow{CD}=\left(\frac{h}{2},-q,\frac{h}{2}\right)\)，\(\overrightarrow{DE}=\left(-\frac{h}{4},0,\frac{h}{4}\right)\). 由此
\[
\overrightarrow{DE}\cdot\overrightarrow{AB_1}=0
\]
以及
\[
\overrightarrow{DE}\cdot\overrightarrow{CD}=0
\]
又因为 \(E\) 在直线 \(AB_1\) 上，\(D\) 在直线 \(CD\) 上，所以 \(DE\) 是异面直线 \(AB_1\) 与 \(CD\) 的公垂线.

\item 由两异面直线所成角的定义及已知条件，得
\[
\frac{\left|\overrightarrow{AB_1}\cdot\overrightarrow{CD}\right|}{\left|\overrightarrow{AB_1}\right|\left|\overrightarrow{CD}\right|}=\cos45^{\circ}=\frac{\sqrt{2}}{2}
\]
代入向量坐标，得
\[
\frac{h^2}{h\sqrt{2}\sqrt{\frac{h^2}{2}+q^2}}=\frac{\sqrt{2}}{2}
\]
所以 \(q^2=\frac{h^2}{2}\)，即 \(q=\frac{h}{\sqrt{2}}\). 取平面 \(A_1AC_1\) 和 \(B_1AC_1\) 的法向量
\[
\bs{n}_1=\overrightarrow{AA_1}\times\overrightarrow{AC_1}=\left(-\frac{h^2}{\sqrt{2}},\frac{h^2}{2},0\right)
\]
以及
\[
\bs{n}_2=\overrightarrow{AB_1}\times\overrightarrow{AC_1}=\left(-\frac{h^2}{\sqrt{2}},-\frac{h^2}{2},\frac{h^2}{\sqrt{2}}\right)
\]
设二面角 \(A_1-AC_1-B_1\) 的大小为 \(\varphi\)，则
\[
\cos\varphi=\frac{\bs{n}_1\cdot\bs{n}_2}{|\bs{n}_1||\bs{n}_2|}=\frac{h^4/4}{(\sqrt{3}h^2/2)(\sqrt{5}h^2/2)}=\frac{1}{\sqrt{15}}
\]
因此二面角 \(A_1-AC_1-B_1\) 的大小为 \(\arccos\frac{1}{\sqrt{15}}\).
\end{enumerate}
\end{solution}

\begin{problem}
如图，由 \(M\) 到 \(N\) 的电路中有 \(4\) 个组件，分别标为 \(T_{1}\)，\(T_{2}\)，\(T_{3}\)，\(T_{4}\)，电流能通过 \(T_{1}\)，\(T_{2}\)，\(T_{3}\) 的概率都是 \(p\)，电流能通过 \(T_{4}\) 的概率是 \(0.9\). 电流能通过各组件相互独立. 已知 \(T_{1}\)，\(T_{2}\)，\(T_{3}\) 中至少有一个能通过电流的概率为 \(0.999\).
\begin{enumerate}
\item 求 \(p\)；
\item 求电流能在 \(M\) 与 \(N\) 之间通过的概率；
\item \(\xi\) 表示 \(T_{1}\)，\(T_{2}\)，\(T_{3}\)，\(T_{4}\) 中能通过电流的组件个数，求 \(\xi\) 的期望.
\end{enumerate}
\bitmapfigure[width=0.60\linewidth]{img/2010/syllabus_paper_2_science/q20_fig1.png}
\end{problem}

\begin{solution}
设事件 \(A_i\) 表示电流能通过组件 \(T_i\).

\begin{enumerate}
\item 由独立性，前三个组件中至少有一个能通过电流的概率为
\[
1-P\left(\overline{A_1}\cap\overline{A_2}\cap\overline{A_3}\right)=1-(1-p)^3=0.999
\]
由于 \(p\) 是概率，\(0\leqslant p\leqslant1\)，所以 \(1-p=0.1\)，即 \(p=0.9\).

\item 由电路图可知，除 \(T_4\) 所在支路外，另一条通路要求 \(T_3\) 能通过电流且 \(T_1\)，\(T_2\) 中至少一个能通过电流. 因此电流能在 \(M\) 与 \(N\) 之间通过的概率为
\[
P(A_4)+P(\overline{A_4})P\bigl(A_3\cap(A_1\cup A_2)\bigr)=0.9+0.1p\bigl[1-(1-p)^2\bigr]
\]
代入 \(p=0.9\)，得
\[
0.9+0.1\times0.9\times(1-0.1^2)=0.9891
\]

\item 令 \(I_i\) 为事件 \(A_i\) 的示性变量，则 \(\xi=I_1+I_2+I_3+I_4\). 由期望的线性性质，
\[
E(\xi)=E(I_1)+E(I_2)+E(I_3)+E(I_4)=3p+0.9=3.6
\]
\end{enumerate}
\end{solution}

\begin{problem}
已知斜率为 \(1\) 的直线 \(l\) 与双曲线 \(C: \frac{x^{2}}{a^{2}} - \frac{y^{2}}{b^{2}} = 1\)（\(a > 0\)，\(b > 0\)）相交于 \(B\)，\(D\) 两点，且 \(BD\) 的中点为 \(M(1,3)\).
\begin{enumerate}
\item 求 \(C\) 的离心率；
\item 设 \(C\) 的右顶点为 \(A\)，右焦点为 \(F\)，\(|DF| \cdot |BF| = 17\)，证明：过 \(A\)，\(B\)，\(D\) 三点的圆与 \(x\) 轴相切.
\end{enumerate}
\end{problem}

\begin{solution}
\begin{enumerate}
\item 由直线 \(l\) 的斜率及其中点 \(M(1,3)\)，得直线方程为 \(y=x+2\). 设 \(B\)，\(D\) 的横坐标分别为 \(x_1\)，\(x_2\)，将直线方程代入双曲线方程，得
\[
\frac{x^2}{a^2}-\frac{(x+2)^2}{b^2}=1
\]
整理为关于 \(x\) 的二次方程，由韦达定理及 \(x_1+x_2=2\)，得
\[
x_1+x_2=\frac{4a^2}{b^2-a^2}=2
\]
所以 \(b^2=3a^2\)，双曲线的离心率为
\[
e=\sqrt{1+\frac{b^2}{a^2}}=2
\]
因此 \(C\) 的离心率为 \(2\).

\item
由 \(c^2=a^2+b^2=4a^2\)，可知右焦点为 \(F=(2a,0)\). 令 \(x_1=1-r\)，\(x_2=1+r\)，其中 \(r>0\). 将 \(b^2=3a^2\) 代回交点方程，得
\[
3x^2-(x+2)^2=3a^2
\]
即
\[
2x^2-4x-4=3a^2
\]
从而 \(r^2=3+\frac{3a^2}{2}\). 对应两交点的纵坐标分别为 \(3-r\)，\(3+r\)，于是
\[
|BF|^2=(1-r-2a)^2+(3-r)^2
\]
以及
\[
|DF|^2=(1+r-2a)^2+(3+r)^2
\]
记 \(K=7a^2-4a+16\)，则由上式及 \(r^2=3+\frac{3a^2}{2}\)，有
\[
|BF|^2=K-2r(4-2a)
\]
以及
\[
|DF|^2=K+2r(4-2a)
\]
因此
\[
\begin{gathered}
|BF|^2|DF|^2
 =K^2-4r^2(4-2a)^2\\
 =(7a^2-4a+16)^2-4\left(3+\frac{3a^2}{2}\right)(4-2a)^2\\
 =49a^4-56a^3+240a^2-128a+256\\
 \mathrel{\phantom{=}}-(24a^4-96a^3+144a^2-192a+192)\\
 =(5a^2+4a+8)^2
\end{gathered}
\]
两距离均为正，结合 \(|DF|\cdot|BF|=17\)，得
\[
5a^2+4a+8=17
\]
即
\[
5a^2+4a-9=0
\]
所以 \((5a+9)(a-1)=0\). 由于 \(a>0\)，故 \(a=1\). 此时交点横坐标满足
\[
2x^2-4x-7=0
\]
设圆 \(\Gamma\) 的方程为
\[
x^2+y^2-2x-6y+1=0
\]
点 \(A=(1,0)\) 满足该方程. 对直线 \(l\) 上任意点 \(P=(x,x+2)\)，代入圆方程左端得
\[
x^2+(x+2)^2-2x-6(x+2)+1=2x^2-4x-7
\]
所以当 \(x\) 为交点横坐标时，\(B\)，\(D\) 也在圆 \(\Gamma\) 上. 又
\[
x^2+y^2-2x-6y+1=0\quad\Longleftrightarrow\quad (x-1)^2+(y-3)^2=9
\]
该圆心为 \((1,3)\)，半径为 \(3\)，与 \(x\) 轴的交点满足 \((x-1)^2=0\)，只有 \(A=(1,0)\). 因此过 \(A\)，\(B\)，\(D\) 三点的圆与 \(x\) 轴相切.
\end{enumerate}
\end{solution}

\section{选考题：解答题}

\begin{problem}
设函数 \( f(x) = 1 - \e^{-x} \).
\begin{enumerate}
\item 证明：当 \(x > -1\) 时，\(f(x) \geqslant \frac{x}{x + 1}\)；
\item 设当 \(x \geqslant 0\) 时，\(f(x) \leqslant \frac{x}{ax + 1}\)，求 \(a\) 的取值范围.
\end{enumerate}
\end{problem}

\begin{solution}
\begin{enumerate}
\item 当 \(x>-1\) 时，\(x+1>0\)，且
\[
f(x)-\frac{x}{x+1}=1-\e^{-x}-\frac{x}{x+1}=\frac{\e^x-x-1}{(x+1)\e^x}
\]
令 \(g(x)=\e^x-x-1\)，则 \(g'(x)=\e^x-1\). 当 \(x<0\) 时，\(g'(x)<0\)；当 \(x>0\) 时，\(g'(x)>0\)，所以 \(g(x)\) 在 \(x=0\) 处取得最小值 \(g(0)=0\). 故 \(f(x)\geqslant\frac{x}{x+1}\).

\item 若 \(a<0\)，则 \(x=-1/a>0\) 时 \(ax+1=0\)，右端无定义；并且取 \(x>-1/a\) 时右端为负，而 \(f(x)=1-\e^{-x}>0\)，不等式也不可能对所有 \(x\geqslant0\) 成立. 因此必须有 \(a\geqslant0\). 此时 \(ax+1>0\)，原不等式等价于
\[
a\leqslant\frac{x-1+\e^{-x}}{x(1-\e^{-x})}=:\varphi(x)
\]
又
\[
\varphi(x)-\frac{1}{2}=\frac{x-2+(x+2)\e^{-x}}{2x(1-\e^{-x})}
\]
令 \(h(x)=x-2+(x+2)\e^{-x}\)，则 \(h(0)=0\)，且
\[
h'(x)=1-(x+1)\e^{-x}=\e^{-x}\bigl(\e^x-x-1\bigr)\geqslant0
\]
所以 \(h(x)\geqslant0\)，从而 \(\varphi(x)\geqslant\frac{1}{2}\). 另一方面，分子、分母在 \(x\to0^+\) 时都趋于 \(0\)，连续使用两次洛必达法则得
\[
\lim_{x\to0^+}\varphi(x)
=\lim_{x\to0^+}\frac{1-\e^{-x}}{1-\e^{-x}+x\e^{-x}}
=\lim_{x\to0^+}\frac{\e^{-x}}{(2-x)\e^{-x}}
=\frac{1}{2}
\]
故必要条件为 \(0\leqslant a\leqslant\frac{1}{2}\). 当 \(0\leqslant a\leqslant\frac{1}{2}\) 时，由 \(a\leqslant\varphi(x)\) 可反推得原不等式对一切 \(x>0\) 成立，且 \(x=0\) 时两边均为 \(0\). 因此 \(a\) 的取值范围是 \(\left[0,\frac{1}{2}\right]\).
\end{enumerate}
\end{solution}
